\documentclass[11pt]{article}
\usepackage{../../../tex/lecturestyle}

\title{Zeta \& Mobius Transform (Subset Relations)}
\author{}
\date{}

\begin{document}
\thispagestyle{fancy}

\lecturetitleband{Set Operations \& Functions}{Zeta \& M\"obius Transform (Subset Relations)}{4}

\section{Overview}
This lecture covers the zeta and M\"obius transforms with respect to the
subset and superset relations on bitmasks. The subset zeta transform is
also widely known in competitive programming as ``sum over subsets''
(SOS): given $A_0,\ldots,A_{2^n-1}$ indexed by subsets of
$\{0,\ldots,n-1\}$, compute for every $T$ the sum of $A_S$ over all
$S \subseteq T$. Naively this is $O(4^n)$; a submask-enumeration trick
brings it to $O(3^n)$; and a bit-by-bit technique --- the \emph{fast
zeta transform} --- computes it (and its inverse, the \emph{fast M\"obius
transform}) in $O(n\cdot 2^n)$. As an application, this lecture derives
$O(n\cdot 2^n)$ algorithms for OR-convolution and AND-convolution of two
length-$2^n$ arrays. Zeta/M\"obius transforms and these bitwise
convolutions are more generally defined for arbitrary partially ordered
sets and lattices; this lecture treats only the subset/superset case.
Assumes representing subsets of $\{0,\ldots,n-1\}$ as bitmasks.

\section{The Subset Zeta Transform}
\begin{definitionbox}
For $A=(A_0,\ldots,A_{2^n-1})$, the sequence $B=(B_0,\ldots,B_{2^n-1})$
defined by
\[
  B_T = \sum_{S \subseteq T} A_S
\]
is the \emph{subset zeta transform} of $A$ (also called SOS, ``sum over
subsets'').
\end{definitionbox}

\begin{examplebox}
For $A=(3,1,4,1,5,9,2,6)$ (indices $0$ through $7$, i.e.\ $n=3$):
$B_0=A_0=3$; $B_1=A_0+A_1=4$; $B_3=A_0+A_1+A_2+A_3=9$; and so on, giving
$B=(3,4,7,9,8,18,14,31)$.
\end{examplebox}

\subsection{An \texorpdfstring{$O(4^n)$}{O(4\^{}n)} Algorithm}
The definition itself gives an $O(4^n)$ algorithm: for every pair
$(S,T)$ with $S \subseteq \{0,\ldots,n-1\}$, $T \subseteq
\{0,\ldots,n-1\}$, test whether $S \subseteq T$ (as bitmasks, whether
\texttt{(S \& T) == S}) and if so add $A_S$ into $B_T$. This costs
$O(4^n)$ since there are $2^n$ choices each for $S$ and $T$.

\subsection{An \texorpdfstring{$O(3^n)$}{O(3\^{}n)} Algorithm}
Rather than looping over all $S$ and testing $S \subseteq T$, looping
only over the submasks of $T$ directly (via the standard
\texttt{for (s = t; ; s = (s - 1) \& t)} idiom, stopping after
processing $s=0$) computes the same sum in $O(3^n)$ overall.

\begin{theorembox}
The number of pairs $(S,T)$ with $S \subseteq T \subseteq
\{0,\ldots,n-1\}$ is $3^n$.
\end{theorembox}
\begin{proof}
Building such a pair amounts to choosing, independently for each
$i=0,\ldots,n-1$, one of three options: $i \notin S, i \notin T$;
$i \notin S, i \in T$; or $i \in S, i \in T$ (the fourth combination,
$i \in S, i \notin T$, is excluded by $S \subseteq T$). This gives $3^n$
total choices.
\end{proof}

\subsection{An \texorpdfstring{$O(n\cdot 2^n)$}{O(n*2\^{}n)} Algorithm (Fast Zeta Transform)}
Write $T$ and $S$ in binary as $T=t_{n-1}\cdots t_1t_0$,
$S=s_{n-1}\cdots s_1s_0$. Replace each bit $t_i=1$ of $T$ by a wildcard
$*$, one at a time, to get a sequence of pattern strings over
$\{0,1,*\}$; a bit string $S$ \emph{matches} a pattern if it agrees with
every fixed ($0$ or $1$) position and is unconstrained at every $*$.
Starting from the all-fixed pattern (matching only $S=T$) and replacing
$T$'s $1$-bits with $*$ one at a time, the set of strings matching the
pattern grows to be exactly $\{S : S \subseteq T\}$ once every $1$-bit
has become a wildcard --- so accumulating sums as each bit turns into a
wildcard computes $B_T$ for every $T$ simultaneously.

\begin{theorembox}[Fast Zeta Transform]
For $i=0,\ldots,n-1$ in order: for every $T$ with bit $i$ set, add
$A_{T \oplus 2^i}$ (i.e.\ $A_T$ with bit $i$ cleared) into $A_T$
in place. After all $n$ passes, $A_T$ holds $B_T$. This runs in
$O(n\cdot 2^n)$ time.
\end{theorembox}

\begin{notebox}
For better constant-factor performance, iterate only over the pairs
$(i,T)$ that require work, rather than testing a condition inside a full
double loop: partition $[0,2^n)$ into blocks of size $2^{i+1}$, and for
each block add the lower half (bit $i$ clear) into the upper half (bit
$i$ set) elementwise. This is the same computation, restructured to
avoid a per-iteration branch.
\end{notebox}

% FIGURE: zeta_passes
\begin{figure}[H]
\centering
\begin{tikzpicture}[x=1.45cm,y=1.35cm]
\node[llabel] at (0,0.7) {$000$};
\node[llabel] at (1,0.7) {$001$};
\node[llabel] at (2,0.7) {$010$};
\node[llabel] at (3,0.7) {$011$};
\node[llabel] at (4,0.7) {$100$};
\node[llabel] at (5,0.7) {$101$};
\node[llabel] at (6,0.7) {$110$};
\node[llabel] at (7,0.7) {$111$};
\node[llabel,anchor=east] at (-0.65,0) {input $A$};
\node[lcell,minimum width=0.95cm,] (n0_0) at (0,0) {3};
\node[lcell,minimum width=0.95cm,] (n0_1) at (1,0) {1};
\node[lcell,minimum width=0.95cm,] (n0_2) at (2,0) {4};
\node[lcell,minimum width=0.95cm,] (n0_3) at (3,0) {1};
\node[lcell,minimum width=0.95cm,] (n0_4) at (4,0) {5};
\node[lcell,minimum width=0.95cm,] (n0_5) at (5,0) {9};
\node[lcell,minimum width=0.95cm,] (n0_6) at (6,0) {2};
\node[lcell,minimum width=0.95cm,] (n0_7) at (7,0) {6};
\node[llabel,anchor=east] at (-0.65,-1) {after bit 0};
\node[lcell,minimum width=0.95cm,] (n1_0) at (0,-1) {3};
\node[lcell,minimum width=0.95cm,lhi] (n1_1) at (1,-1) {4};
\node[lcell,minimum width=0.95cm,] (n1_2) at (2,-1) {4};
\node[lcell,minimum width=0.95cm,lhi] (n1_3) at (3,-1) {5};
\node[lcell,minimum width=0.95cm,] (n1_4) at (4,-1) {5};
\node[lcell,minimum width=0.95cm,lhi] (n1_5) at (5,-1) {14};
\node[lcell,minimum width=0.95cm,] (n1_6) at (6,-1) {2};
\node[lcell,minimum width=0.95cm,lhi] (n1_7) at (7,-1) {8};
\node[llabel,anchor=east] at (-0.65,-2) {after bit 1};
\node[lcell,minimum width=0.95cm,] (n2_0) at (0,-2) {3};
\node[lcell,minimum width=0.95cm,] (n2_1) at (1,-2) {4};
\node[lcell,minimum width=0.95cm,lhi] (n2_2) at (2,-2) {7};
\node[lcell,minimum width=0.95cm,lhi] (n2_3) at (3,-2) {9};
\node[lcell,minimum width=0.95cm,] (n2_4) at (4,-2) {5};
\node[lcell,minimum width=0.95cm,] (n2_5) at (5,-2) {14};
\node[lcell,minimum width=0.95cm,lhi] (n2_6) at (6,-2) {7};
\node[lcell,minimum width=0.95cm,lhi] (n2_7) at (7,-2) {22};
\node[llabel,anchor=east] at (-0.65,-3) {after bit 2 $=B$};
\node[lcell,minimum width=0.95cm,] (n3_0) at (0,-3) {3};
\node[lcell,minimum width=0.95cm,] (n3_1) at (1,-3) {4};
\node[lcell,minimum width=0.95cm,] (n3_2) at (2,-3) {7};
\node[lcell,minimum width=0.95cm,] (n3_3) at (3,-3) {9};
\node[lcell,minimum width=0.95cm,lhi] (n3_4) at (4,-3) {8};
\node[lcell,minimum width=0.95cm,lhi] (n3_5) at (5,-3) {18};
\node[lcell,minimum width=0.95cm,lhi] (n3_6) at (6,-3) {14};
\node[lcell,minimum width=0.95cm,lhi] (n3_7) at (7,-3) {31};
\draw[lthin,lectureink!35] (n0_0) -- (n1_0);
\draw[lthin,lectureink!35] (n0_1) -- (n1_1);
\draw[lthin,lectureaccent,thick] (n0_0) -- (n1_1);
\draw[lthin,lectureink!35] (n0_2) -- (n1_2);
\draw[lthin,lectureink!35] (n0_3) -- (n1_3);
\draw[lthin,lectureaccent,thick] (n0_2) -- (n1_3);
\draw[lthin,lectureink!35] (n0_4) -- (n1_4);
\draw[lthin,lectureink!35] (n0_5) -- (n1_5);
\draw[lthin,lectureaccent,thick] (n0_4) -- (n1_5);
\draw[lthin,lectureink!35] (n0_6) -- (n1_6);
\draw[lthin,lectureink!35] (n0_7) -- (n1_7);
\draw[lthin,lectureaccent,thick] (n0_6) -- (n1_7);
\draw[lthin,lectureink!35] (n1_0) -- (n2_0);
\draw[lthin,lectureink!35] (n1_1) -- (n2_1);
\draw[lthin,lectureink!35] (n1_2) -- (n2_2);
\draw[lthin,lectureaccent,thick] (n1_0) -- (n2_2);
\draw[lthin,lectureink!35] (n1_3) -- (n2_3);
\draw[lthin,lectureaccent,thick] (n1_1) -- (n2_3);
\draw[lthin,lectureink!35] (n1_4) -- (n2_4);
\draw[lthin,lectureink!35] (n1_5) -- (n2_5);
\draw[lthin,lectureink!35] (n1_6) -- (n2_6);
\draw[lthin,lectureaccent,thick] (n1_4) -- (n2_6);
\draw[lthin,lectureink!35] (n1_7) -- (n2_7);
\draw[lthin,lectureaccent,thick] (n1_5) -- (n2_7);
\draw[lthin,lectureink!35] (n2_0) -- (n3_0);
\draw[lthin,lectureink!35] (n2_1) -- (n3_1);
\draw[lthin,lectureink!35] (n2_2) -- (n3_2);
\draw[lthin,lectureink!35] (n2_3) -- (n3_3);
\draw[lthin,lectureink!35] (n2_4) -- (n3_4);
\draw[lthin,lectureaccent,thick] (n2_0) -- (n3_4);
\draw[lthin,lectureink!35] (n2_5) -- (n3_5);
\draw[lthin,lectureaccent,thick] (n2_1) -- (n3_5);
\draw[lthin,lectureink!35] (n2_6) -- (n3_6);
\draw[lthin,lectureaccent,thick] (n2_2) -- (n3_6);
\draw[lthin,lectureink!35] (n2_7) -- (n3_7);
\draw[lthin,lectureaccent,thick] (n2_3) -- (n3_7);
\end{tikzpicture}
\caption{The fast zeta transform on $A=(3,1,4,1,5,9,2,6)$ with $n=3$. Pass $i$ adds, to every entry whose index has bit $i$ set, the entry whose index differs in that bit (blue arrows); shaded cells are the ones that change. After the three passes each entry is the sum over all subsets of its index, e.g.\ $B_7=31=\sum_S A_S$.}\label{fig:zeta_passes}
\end{figure}

\section{The Subset M\"obius Transform}
\begin{definitionbox}
Given $B=(B_0,\ldots,B_{2^n-1})$, the sequence
$A=(A_0,\ldots,A_{2^n-1})$ satisfying $B_T = \sum_{S\subseteq T}A_S$ for
every $T$ is the \emph{subset M\"obius transform} of $B$ --- the inverse
of the subset zeta transform. (Inclusion--exclusion gives the closed
form $A_S = \sum_{T\subseteq S}(-1)^{|S|-|T|}B_T$, though this lecture
does not need it.)
\end{definitionbox}

\begin{theorembox}[Fast M\"obius Transform]
Running the fast zeta transform's bit-flip step in reverse --- for
$i=0,\ldots,n-1$, for every $T$ with bit $i$ set, \emph{subtract}
$A_{T\oplus 2^i}$ from $A_T$ --- inverts the fast zeta transform exactly,
in $O(n\cdot 2^n)$ time. (The order in which the $n$ bit positions are
processed does not matter for either the fast zeta or fast M\"obius
transform, since each pass only touches, for each pair differing in bit
$i$, that bit's own contribution.)
\end{theorembox}

\section{Superset Zeta and M\"obius Transforms}
\begin{definitionbox}
For $A=(A_0,\ldots,A_{2^n-1})$, the \emph{superset zeta transform} is
$B_T = \sum_{S \supseteq T}A_S$, and its inverse is the \emph{superset
M\"obius transform}.
\end{definitionbox}
The algorithms are the mirror image of the subset case: in the pattern
matching view, $T$'s $0$-bits (rather than $1$-bits) are replaced by
wildcards one at a time. Both run in $O(n\cdot 2^n)$.

\section{OR-Convolution and AND-Convolution}
A less immediately obvious application of these transforms is computing
convolutions with respect to bitwise OR and AND.

\begin{definitionbox}[OR-convolution]
Given $A,B$ indexed by subsets of $\{0,\ldots,n-1\}$, compute
$C_U = \sum_{S\cup T=U} A_SB_T$.
\end{definitionbox}
The condition $S\cup T=U$ is awkward to work with directly, but
$S\cup T\subseteq U$ factors nicely: $S\cup T\subseteq U \iff
S\subseteq U \text{ and } T\subseteq U$. So define
\[
  D_U = \sum_{S\cup T\subseteq U} A_SB_T
  = \Bigl(\sum_{S\subseteq U}A_S\Bigr)\Bigl(\sum_{T\subseteq U}B_T\Bigr),
\]
computable as a pointwise product of subset zeta transforms. Relating
$D$ back to $C$:
\[
  D_U = \sum_{U'\subseteq U} C_{U'},
\]
i.e.\ $D$ is the subset zeta transform of $C$ --- equivalently, $C$ is
the subset M\"obius transform of $D$.

\begin{theorembox}
OR-convolution: take the subset zeta transform of $A$ and of $B$,
multiply pointwise, then take the subset M\"obius transform of the
result. This computes $C$ in $O(n\cdot 2^n)$.
\end{theorembox}

\begin{definitionbox}[AND-convolution]
Given $A,B$ indexed by subsets of $\{0,\ldots,n-1\}$, compute
$C_U = \sum_{S\cap T=U} A_SB_T$.
\end{definitionbox}
\begin{theorembox}
AND-convolution: take the superset zeta transform of $A$ and of $B$,
multiply pointwise, then take the superset M\"obius transform of the
result (the argument mirrors the OR-convolution case, using
$S\cap T\supseteq U \iff S\supseteq U \text{ and } T \supseteq U$). This
also runs in $O(n\cdot 2^n)$.
\end{theorembox}

\section{Generalizing the Operation}
The fast zeta transform generalizes to any commutative monoid in place
of integer addition; the fast M\"obius transform (which needs to
subtract) generalizes to any commutative group; and OR-/AND-convolution
(which needs both addition and multiplication) generalize to any ring.

\section{Reference Implementations}
The examples below implement subset zeta/M\"obius transforms three ways
each (an $O(4^n)$ definition-level double loop, an $O(3^n)$
submask-enumeration loop, and the $O(n\cdot 2^n)$ fast zeta/M\"obius
transform), reproduce the lecture's own worked example exactly, and
cross-check the three implementations against each other and the
M\"obius transform against being an exact inverse of the zeta transform
--- for both the subset and superset relations --- over many random
trials. OR- and AND-convolution are then implemented via the
zeta/pointwise-product/M\"obius pipeline and cross-checked against an
$O(4^n)$ brute-force convolution definition.

\codenote
\subsection{Rust}
\corelisting{Rust}{example.rs}

\subsection{Python}
\corelisting{Python}{example.py}

\section{Worked Example}
\begin{examplebox}
For $A=(3,1,4,1,5,9,2,6)$, all three subset zeta transform
implementations reproduce the lecture's own
$B=(3,4,7,9,8,18,14,31)$ exactly. Across 100 random trials each (with
$n$ up to $8$), the $O(4^n)$, $O(3^n)$, and $O(n\cdot 2^n)$
implementations of both the subset and superset zeta transforms agree
exactly, and applying the corresponding M\"obius transform exactly
recovers the original array in every case. OR- and AND-convolution
(via the zeta/pointwise-product/M\"obius pipeline) match an $O(4^n)$
brute-force convolution over 100 further random trials each.
\end{examplebox}

\section{Further Reading}
\begin{itemize}[leftmargin=1.5em]
  \item AtCoder Algorithm Lectures index (source of this curriculum):
        \url{https://info.atcoder.jp/entry/algorithm_lectures/index}
  \item Library Checker, ``Bitwise And Convolution'':
        \url{https://judge.yosupo.jp/problem/bitwise_and_convolution}
  \item AtCoder ARC100 E, ``Or Plus Max'':
        \url{https://atcoder.jp/contests/arc100/tasks/arc100_c}
  \item USACO Guide, ``Sum over Subsets DP'':
        \url{https://usaco.guide/plat/dp-sos}
\end{itemize}

\end{document}
